\section{Instruksi Mesin, Direktif Assembler, dan Direktif Data}
\label{sec:instruksi-direktif}

Satu berkas kode sumber assembly tidak hanya berisi instruksi.
Sebagian besar baris justru berupa \term{direktif}, yaitu perintah
yang ditujukan kepada assembler, bukan kepada prosesor. Bab ini
membuka perbedaan keduanya, lalu membahas direktif data.

\subsection{Dua Jenis Baris dalam Kode Sumber}

Setiap baris berkas \file{.asm} dapat digolongkan menjadi salah satu
dari dua kelompok berikut.

\begin{description}
    \item[Instruksi mesin] Baris yang mnemonic-nya (misalnya
          \asm{MOV}, \asm{ADD}, \asm{INT}) akan diterjemahkan
          menjadi \term{opcode}, yaitu kode mesin yang benar-benar
          dieksekusi prosesor. Instruksi hanya bekerja pada saat
          program berjalan (\eng{run time}). Bentuk lengkap sebuah
          instruksi beserta operandnya telah dibahas pada
          Bab~\ref{chap:format-instruksi}.

    \item[Direktif assembler] Baris yang hanya berlaku pada saat
          perakitan (\eng{assemble time}). Direktif memberi tahu
          assembler bagaimana menata program: di mana data harus
          diletakkan, berapa besar ruang yang disediakan, nilai apa
          yang dipakai sebagai konstanta, dan seterusnya. Direktif
          disebut juga \eng{pseudo-op} karena bentuknya mirip
          instruksi tetapi \important{tidak menghasilkan kode mesin
          apa pun}.
\end{description}

Perbedaan ini paling mudah dilihat pada berkas \file{.lst}. Di
sana, baris instruksi memiliki kolom kode mesin, sedangkan baris
direktif tidak.

\begin{session}[caption={Cuplikan \file{.lst}: direktif tidak memiliki kode mesin.}, label={lst:lst-direktif}]
  4                             .DATA
  5  0000  48 61 6C 6F        pesan   DB 'Halo'
  6                             panjang EQU 5
  7
  8                             .CODE
  9  0000  B8 ----            START:  MOV AX, @DATA
 10  0003  8E D8                      MOV DS, AX
\end{session}

Baris 6 berisi direktif \dir{EQU}. Kolom kode mesinnya kosong karena
tidak ada instruksi yang dihasilkan, sedangkan baris 9 dan 10
menghasilkan opcode \hex{B8} dan \hex{8E D8}. Tanda \code{----}
menandai nilai yang baru diselesaikan linker
(Bagian~\ref{sec:peran-assembler-linker}).

\begin{infobox}
Direktif yang mengisi memori (\dir{DB}, \dir{DW}, \dir{DD}) tetap
mempengaruhi isi berkas \file{.exe}. Yang tidak dihasilkannya adalah
\emph{instruksi}. Karena itu, pertanyaan ``apakah direktif
menghasilkan kode mesin?'' harus dijawab dengan membedakan kode
mesin (\eng{opcode}) dan data.
\end{infobox}

\subsection{Direktif Data}

Direktif data menyatakan isi memori, sedangkan direktif lain mengatur
nama, letak, dan struktur program. Tabel~\ref{tab:direktif-lengkap}
merangkum seluruh direktif yang dibahas pada bab ini.

\begin{table}[H]
    \centering
    \caption{Ringkasan direktif assembler 8086 beserta fungsinya.}
    \label{tab:direktif-lengkap}
    \small
    \begin{tabular}{@{}p{1.9cm}p{4.1cm}p{3.4cm}p{4.6cm}@{}}
        \toprule
        \textbf{Direktif} & \textbf{Fungsi} & \textbf{Contoh} &
        \textbf{Ukuran / efek} \\
        \midrule
        \dir{DB} & Data 1 byte &
        \lstasm|DB 0Ah| & 1 byte per nilai \\
        \addlinespace
        \dir{DW} & Data 2 byte &
        \lstasm|DW 1234h| & 2 byte; \eng{little-endian} \\
        \addlinespace
        \dir{DD} & Data 4 byte &
        \lstasm|DD 12345678h| & 4 byte; \eng{little-endian} \\
        \addlinespace
        \dir{DT} & Data 10 byte &
        \lstasm|DT 1.0| & 10 byte (bilangan pecahan) \\
        \addlinespace
        \dir{DUP} & Mengulang satu pola &
        \lstasm|DB 10 DUP(0)| & 10 byte bernilai nol \\
        \addlinespace
        \code{?} & Ruang tanpa nilai awal &
        \lstasm|DB 8 DUP(?)| & 8 byte tanpa isi \\
        \addlinespace
        \code{\$} & Penghitung lokasi &
        \lstasm|EQU $ - pesan| & Nilai \eng{offset} saat itu \\
        \addlinespace
        \dir{EQU} & Konstanta simbolik tetap &
        \lstasm|KELUAR EQU 4Ch| & 0 byte; tidak dapat diubah \\
        \addlinespace
        \code{=} & Simbol yang boleh berubah &
        \lstasm|NILAI = 10| & 0 byte; boleh didefinisikan ulang \\
        \addlinespace
        \dir{ORG} & Menetapkan penghitung lokasi &
        \lstasm|ORG 100h| & 0 byte; menggeser \eng{offset} \\
        \addlinespace
        \dir{SEGMENT} & Membuka segmen bernama &
        \lstasm|kode SEGMENT| & 0 byte; mengawali segmen \\
        \addlinespace
        \dir{ENDS} & Menutup segmen &
        \lstasm|kode ENDS| & 0 byte; wajib bagi setiap
        \dir{SEGMENT} \\
        \addlinespace
        \dir{ASSUME} & Menetapkan segmen \eng{default} &
        \lstasm|ASSUME CS:kode| & 0 byte; hanya untuk assembler \\
        \addlinespace
        \dir{END} & Menutup berkas sumber &
        \lstasm|END START| & 0 byte; menetapkan titik masuk \\
        \addlinespace
        \dir{.MODEL} & Memilih model memori &
        \lstasm|.MODEL SMALL| & 0 byte; menentukan jumlah segmen \\
        \addlinespace
        \dir{PROC}, \dir{ENDP} & Membuka dan menutup prosedur &
        \lstasm|cetak PROC| & 0 byte; menandai batas prosedur \\
        \bottomrule
    \end{tabular}
\end{table}

Beberapa bentuk penulisan berikut dipakai terus-menerus pada
program nyata.

\begin{enumerate}
    \item \textbf{Literal string.} \lstasm|DB 'Halo'| menyediakan
          empat byte: \hex{48}, \hex{61}, \hex{6C}, \hex{6F}.
          String tidak otomatis diakhiri apa pun, sehingga penanda
          akhir harus ditambahkan sendiri, misalnya
          \lstasm|DB 'Halo', '$'| untuk layanan DOS fungsi
          \hex{09}.

    \item \textbf{Beberapa nilai sekaligus.}
          \lstasm|DW 1, 2, 3| menyediakan tiga kata, yaitu enam
          byte, dengan urutan sesuai penulisan.

    \item \textbf{Operator \dir{DUP}.} Bentuk
          \code{<jumlah> DUP(<nilai>)} mengulang satu pola:
          \lstasm|DB 10 DUP(0)| menyediakan sepuluh byte bernilai
          nol. Bentuk bersarang juga diizinkan, misalnya
          \lstasm|DB 2 DUP(3 DUP(0))|.

    \item \textbf{Penanda \code{?}.} Tanda tanya berarti sediakan
          ruang tanpa nilai awal, seperti pada
          \lstasm|buffer DB 8 DUP(?)|.

    \item \textbf{Penghitung lokasi \code{\$}.} Simbol \code{\$}
          berisi \eng{offset} baris yang sedang dirakit, sehingga
          \lstasm|panjang EQU $ - pesan| menghitung panjang string
          secara otomatis.
\end{enumerate}

\subsection{Urutan \eng{Little-Endian} pada \dir{DW} dan \dir{DD}}

Prosesor 8086 menyimpan bilangan multi-byte dengan urutan
\term{little-endian}: \important{byte berpangkat rendah disimpan pada
alamat yang lebih kecil}. Jadi, \lstasm|nilai DW 1234h|
\emph{tidak} tersimpan sebagai \hex{12} lalu \hex{34}, melainkan
\hex{34} lalu \hex{12}; aturan yang sama berlaku untuk \dir{DD}.

\begin{figure}[H]
    \centering
    \begin{tikzpicture}[
        sel/.style={
            draw=obeBlue, line width=0.8pt,
            minimum width=1.25cm, minimum height=0.8cm,
            align=center, font=\ttfamily\small,
            fill=obeBlueLight, inner sep=0pt
        },
        selD/.style={
            draw=obeTeal, line width=0.8pt,
            minimum width=1.25cm, minimum height=0.8cm,
            align=center, font=\ttfamily\small,
            fill=obeTealLight, inner sep=0pt
        },
        addr/.style={font=\scriptsize, text=obeGray},
        nama/.style={font=\small\bfseries, text=obeBlue, anchor=east},
        catat/.style={font=\scriptsize\itshape, text=obeMaroon}
    ]
        % ---- DW 1234h : dua byte ----
        \node[nama] at (-0.2,0) {\dir{DW} \hex{1234}};
        \node[sel] (w0) at (0,0) {34};
        \node[sel] (w1) at (1.25,0) {12};
        \node[addr, below=0.08cm of w0] {\hex{0300}};
        \node[addr, below=0.08cm of w1] {\hex{0301}};
        \node[catat, above=0.15cm of w1, xshift=-0.62cm]
            {byte rendah lebih dulu};

        % ---- DD 12345678h : empat byte ----
        \node[nama] at (-0.2,-2.3) {\dir{DD} \hex{12345678}};
        \node[selD] (d0) at (0,-2.3) {78};
        \node[selD] (d1) at (1.25,-2.3) {56};
        \node[selD] (d2) at (2.50,-2.3) {34};
        \node[selD] (d3) at (3.75,-2.3) {12};
        \node[addr, below=0.08cm of d0] {\hex{0302}};
        \node[addr, below=0.08cm of d1] {\hex{0303}};
        \node[addr, below=0.08cm of d2] {\hex{0304}};
        \node[addr, below=0.08cm of d3] {\hex{0305}};
        \node[catat, above=0.15cm of d1, xshift=0.62cm]
            {urutan penulisan nilai tetap dibaca dari kiri ke kanan};

        \draw[-{Stealth[length=2mm]}, obeGray, line width=0.8pt]
            (2.9,-0.55) -- (2.9,-1.75);
        \node[font=\scriptsize, text=obeGray, right] at (2.95,-1.15)
            {alamat bertambah};
    \end{tikzpicture}
    \caption{Penataan byte \dir{DW} dan \dir{DD} di memori dengan
    urutan \eng{little-endian}.}
    \label{fig:little-endian}
\end{figure}

Gambar~\ref{fig:little-endian} menjelaskan mengapa urutan ini
penting bagi penyusunan tabel dan \eng{array}. Bila program membaca
satu kata sekaligus dengan \lstasm|MOV AX, nilai|, prosesor sendiri
yang menggabungkan kedua byte dalam urutan yang benar. Sebaliknya,
bila memori dibaca \emph{byte per byte}, urutan little-endian
terlihat dan harus diperhitungkan, seperti dimanfaatkan pada
Kode~\ref{lst:direktif-data}.
